% Generated draft: requires mathematical and source review.
\documentclass[11pt]{article}
\usepackage[margin=25mm]{geometry}
\usepackage{fontspec}
\setmainfont{Latin Modern Roman}
\usepackage{amsmath,amssymb,mathtools,unicode-math}
\setmathfont{Latin Modern Math}
\usepackage{enumitem,needspace,xurl,hyperref}
\setlength{\parindent}{0pt}
\setlength{\parskip}{5pt}
\setlength{\emergencystretch}{3em}
\newcommand{\N}{\mathbb{N}}
\newcommand{\Z}{\mathbb{Z}}
\newcommand{\Q}{\mathbb{Q}}
\newcommand{\R}{\mathbb{R}}
\newcommand{\C}{\mathbb{C}}
\newcommand{\F}{\mathbb{F}}
\newcommand{\A}{\mathbb{A}}
\newcommand{\PP}{\mathbb{P}}
\newcommand{\E}{\mathbb{E}}
\newcommand{\eps}{\varepsilon}
\newcommand{\dd}{\,\mathrm{d}}
\begin{document}
\section*{UnsolvedMath ID: 30000854}
\textbf{Source problem code:} catalogue entry
\subsection{Definitions and mathematical statement}
Let $S$ be a minimal complex surface of general type, understood in the standard smooth connected projective setting: $\kappa(S)=2$ and $S$ has no exceptional curve of self-intersection $-1$. Write $K_S$ for its canonical divisor, $K_S^2$ for its self-intersection, and \[\chi(S)=\chi(\mathcal O_S)=\sum_{i=0}^2(-1)^i\dim_\C H^i(S,\mathcal O_S).\] Thus $\chi$ is the holomorphic Euler characteristic, not the topological one. Let $\pi_1^{\mathrm{alg}}(S)$ be the profinite algebraic fundamental group classifying finite \'etale covers. An irregular \'etale cover means a connected finite \'etale cover $T\to S$ with $q(T)=\dim_\C H^1(T,\mathcal O_T)>0$; absence of such covers includes $q(S)=0$.\par Miles Reid's conjecture in the source asks whether \[K_S^2<4\chi(S)\] implies that $\pi_1^{\mathrm{alg}}(S)$ is, up to finite group extensions, the fundamental group of a curve. To compare profinite groups, the curve group must be read as an algebraic fundamental group, or as the profinite completion of the topological fundamental group in the compact complex-curve setting. The phrase up to finite group extensions does not specify the finite-kernel/finite-index diagram. For example, passing to an open subgroup and then quotienting by a finite normal subgroup is one possible convention, but it is not asserted as the report's exact formulation without further evidence. The curve class and allowed genus are also unstated; genus zero would allow the trivial curve group and hence finite cases.\par The second question asks for explicit upper bounds on \[|\pi_1^{\mathrm{alg}}(S)|\] under the stronger inequality $K_S^2<3\chi(S)$ and the assumption that $S$ has no irregular finite \'etale cover. No candidate numerical bound or permitted dependence on the surface's invariants is supplied. The 2026-08-21 triage labels the record open, reports the conjecture proved in the irregular and $K_S^2<3\chi(S)$ cases, and locates the unresolved no-irregular-cover range at $3\chi(S)\leq K_S^2<4\chi(S)$ within the conjecture's strict upper threshold. These are supplied literature qualifications, not a newly verified classification or bound.
\subsection{Short English statement}
For minimal general-type surfaces below the slope-four threshold, ask whether the algebraic fundamental group is a curve group up to finite extensions. Below slope three, with no irregular finite \'etale cover, ask for explicit bounds on its order. The profinite interpretation and finite-extension convention require clarification.
\subsection{Sources}
\begin{itemize}\item Local UnsolvedMath registry, permanent ID 30000854, source problem catalogue entry, \emph{Reid's Fundamental-Group Conjecture for General-Type Surfaces}, proposed year 2007. The original attributes the conjecture to Miles Reid; no separate proposer metadata is supplied.\item Supplied citation: \emph{Komplexe Algebraische Geometrie} (2008), Oberwolfach Reports, DOI 10.4171/owr/2007/47, pp.\ 2791--2850: \url{https://doi.org/10.4171/owr/2007/47}. Supplied literature label: original Oberwolfach report.\item Supplied literature label: Mendes Lopes--Pardini, algebraic fundamental groups, \url{https://arxiv.org/abs/math/0512483}.\end{itemize}
\end{document}
